% ============================================================================
% sec-computation.tex — Computational methodology: exhaustive accounts,
% audits, and Lean trust boundaries. Label prefix: com: (all labels owned
% by this file).
% ============================================================================

\section{Computational methodology: exhaustive accounts, audits, and trust
boundaries}\label{sec:computation}

Every machine-backed statement in this paper is quoted at a measured
strength, and this section documents the apparatus that produces the
measurements. There are four layers. First, the exact-value pipeline behind
the constructive parts (Section~\ref{com:subsec-pipeline}), together with a
self-contained recheck script released with the paper
(Section~\ref{com:subsec-recheck}). Second, the re-verification account of
the conditional chain's machine layer
(Section~\ref{com:subsec-kline}). Third, the Lean development: what is
formalized, and where its trust boundaries lie
(Section~\ref{com:subsec-lean}). Fourth, the key-table ledger behind the
covering identity (Proposition~\ref{prop:clacover}) and the independent
audit it passed (Section~\ref{com:subsec-ledger}). Formalization status is
itemized declaration by declaration in Section~\ref{sec:open}.

\subsection{The exhaustive construction pipeline}\label{com:subsec-pipeline}

All exact values and all decompositions cited in the constructive parts were
produced and re-checked by one pipeline: a bitmask exact dynamic program
over edge subsets computes exact values of $\pn(G)$ (and of the endpoint
budget $\beta(G,x)$ of Definition~\ref{int:def-budget}, with an endpoint
constraint); a backtracking exact-cover solver with longest-path-first
enumeration handles larger instances; and an independent
\texttt{verify\_decomposition} routine asserts that a witness consists of
simple paths exactly partitioning the edge set. Randomized steps use fixed
seeds, and every number is reproducible from the scripts and certificate
files released with the paper. In the per-theorem verification table
(Appendix~\ref{com:app-comp}), ``exact'' means computed by an exact algorithm
(dynamic program or exhaustive cover), ``witness'' means a machine-verified
decomposition, and ``census'' means that all graphs of the stated class were
processed.

Three rows of that table indicate its granularity.

\begin{itemize}[leftmargin=1.8em]
\item The odd-column value $\pn(K_{4,n})=(n+1)/2$
(Proposition~\ref{prop:conk4large}) was re-derived independently for every
$n=3,\dots,31$ by a structured exact search over $K_{4,n}$, recovering
$(n+1)/2$ each time; the closed-form construction was verified for
$k=6,\dots,301$ (i.e.\ $n=13,\dots,603$), with witness lists
machine-verified. The printed small-case witnesses of
Proposition~\ref{prop:conk4small} are machine-verified by the same pipeline
(edge-exact partition, simplicity).
\item The witnesses behind Corollary~\ref{cor:concol6} of the $n_2=6$ column
($K_{7,6}$, $K_{9,6}$, $K_{11,6}$, $K_{13,6}$) were each verified by two
independent implementations (the pipeline verifier and an adjacency-matrix
checker); the ring construction was re-derived for $k=7..31$ from the base
$k_0=3$, with overlapping chains from $k_0=4,5$.
\item The negatives on the five open bases
(Observation~\ref{obs:concostwall}) are \emph{truncations, not exclusions}:
more than $14$ million search nodes (column $n_2=8$) and about $0.9$ million
nodes (column $n_2=10$) were spent over every ledger-feasible deficit
multiset, without a hit and without exhausting a single schedule. The
conjectural status of the five bases is therefore unchanged by these
searches, in either direction.
\end{itemize}

\subsection{A self-contained independent recheck script}\label{com:subsec-recheck}

The repository releases a recheck script, \texttt{checks.py}, written from
scratch during the independent review of the paper: it imports none of the
project's research code, consists of ten check functions built from
assertions only, writes its results to \texttt{check\_results.json}, and
exits nonzero if any assertion fails. Running
\texttt{python3 checks.py} from the repository root executes all ten checks
in about three seconds (wall clock on a laptop-class machine; measured
$3.1$\,s for the release version). The coverage is:

\begin{itemize}[leftmargin=1.8em]
\item \emph{Ring closure} (Theorem~\ref{thm:conringclosure}): all $280$
parameter pairs ($h=1,\dots,14$; per $h$ the $20$ even ring sizes $m$ from
$2h+2$ to $2h+41$), checking path simplicity and that the paths partition
the edge set exactly once.
\item \emph{The $K_{4,n}$ construction} (Proposition~\ref{prop:conk4large}):
all $95$ parameter values $k=6,\dots,100$, asserting $k+1$ paths, simplicity,
and the exact $A\times B$ edge partition with multiplicity $1$.
\item \emph{The odd-clique local engine} (Theorem~\ref{thm:conL0}): $61$ odd
values $j=1,3,\dots,121$, checking that the removed arcs form a Hamilton
path, that the remaining $j$ paths have the prescribed lengths, and that the
partition is exact. Scope: odd $j$ only; even $j$ is not verified by the
script (it is open; Remark~\ref{rem:conchaincomp}).
\item \emph{The non-infinity prime-family arithmetic}
(Theorem~\ref{thm:claprimefam}): $86$ primes $p\equiv3\pmod4$ up to $991$,
checking the sieve's allowed regimes and the perfect-square obstruction. This
is a finite recheck of the elementary steps; it does not substitute for the
paper proof and covers the non-infinity regimes only.
\item \emph{The $K_9@36$ census ledger}: the six recorded regime lines'
configuration counts against their closed combinatorial forms, the histogram
sums, the zero recorded hits, and the total-domain arithmetic
$18{,}385{,}772{,}680=\binom{71}{8}+\binom{36}{8}\cdot2^8$. This is ledger
and artifact arithmetic; the census itself is not rerun.
\item \emph{The non-infinity death table}: $28$ full-death points for
$n\le200$, of which $23$ have $n-1$ prime (matching
Theorem~\ref{thm:claprimefam}) and the composite points are
$\{112,116,124,160,188\}$. Again artifact accounting, not a rerun of the
table's feasibility searches.
\item \emph{The corrected Fourier identity} (Theorem~\ref{thm:clafourier}):
the $K_3$ configuration $\{\infty,1,3\}$ at $k=3$ with $\psi(a)=\omega^a$,
reproducing both the failure of the superseded form (residual $-2-\omega$)
and the vanishing of the corrected form \eqref{eq:clarho}, in exact
$\mathbb{Z}[\omega]$ integer arithmetic.
\item \emph{Block versus global degree}
(Section~\ref{int:subsec-blockvsglobal}): the $K_{1,5}$-with-pendant-leaves
witness, confirming that $\EV(G)\cong K_{1,5}$ has a vertex of E-degree $5$
while every block of $\EV(G)$ is a single edge (maximum degree $1$).
\item \emph{The Lean placeholder scan}: a comment-stripped scan of the Lean
sources for \texttt{sorry}, \texttt{admit} and \texttt{axiom}
declarations, finding exactly five occurrences, listed in
Section~\ref{com:subsec-lean}.
\item \emph{Snapshot integrity}: SHA-256 verification of the archived source
snapshot (20 files plus the pinned artifacts). In the review environment the
same check additionally reported the six post-review-edit live files (the
expected repair edits to the working notes and four section files, not frozen
deliverables); that comparison required the live working tree, which is not
part of the released repository, so the released script reports the field as
empty.
\end{itemize}

One reading rule binds the released result file: a PASS records that the
intended recheck succeeded, and for two of the ten checks success means
\emph{reproducing a recorded error or distinction} (the falsified Fourier
form; the block-versus-global-degree witness). A PASS is an execution
record, not a certificate of global correctness.

\paragraph{Census methodology note ($K_9@36$).} The largest
configuration-level account behind the classification programme is
conducted line by line: for each regime line of $K_9@36$, the
configurations are enumerated in full and summarized by a closed-form
count checked against the enumeration, with a histogram-sum completeness
assertion per line. The nine regime lines sum to
$\binom{71}{8}=10{,}639{,}125{,}640$ (the Vandermonde identity is
machine-checked against the per-line counts); the infinity line,
$\binom{36}{8}\cdot2^8=7{,}746{,}647{,}040$, is covered in four shards of
$1{,}936{,}661{,}760$ each under an exact non-overlapping-partition
assertion; every line has zero hits; and the total domain is
$18{,}385{,}772{,}680$ configurations. The released checks verify the
ledger arithmetic of this account (per-line closed forms, histogram sums,
recorded hits, and the total); the census itself is not rerun by them.
The tier of the verdict --- machine universal-closure candidate, not
independently audited --- is the one recorded in
Section~\ref{subsubsec:firstlayers}, and the layered-arithmetic verdict of
the same section provides a mechanism-level cross-check.

\paragraph{Complete-walk counts on the flicker line.} The GEN existence
programme contributes a counting account of a different kind: exact
enumeration of zero-dodge complete walks on the skeleton template family,
$w_1=r-5$, by depth-first search with node budgets and cap tracking,
cross-checked across independent implementations. On the odd tiers the
exact counts are $158$, $1{,}733$, $11{,}466$, $81{,}660$, $484{,}578$,
$3{,}670{,}846$ for $r=19,21,23,25,27,29$ --- six tiers, never presented
as one closed sequence; the $r=29$ count is exact but
class-bound-conditioned and is quoted only with that caveat. One further
exact single point, $23{,}933{,}468$ at $r=31$, stands apart at
exhaustive grade ($j=0$: $22{,}607{,}250$; $j=1$: $1{,}326{,}218$). At
$r=33$ there is no count: the search volume is a quantified hard wall,
and a segmented long run reports none. The $w_1=r-4$ side is zero for all
$r\ge19$ and the even-$r$ side with $w_1=r-5$ is zero by the parity
theorem --- both theorem-level, with the four theorem-relevant cells
($r\in\{20,22\}$, both $w_1$) genuinely exhausted at zero; the theorem
clears the zero-dodge side only, complete walks with a positive
zero-defect existing at $r=20,22$ in the record and remaining untouched.

\subsection{The re-verification account of the conditional chain's machine
layer}\label{com:subsec-kline}

The conditional chain of Section~\ref{int:subsec-budget} (Theorem~\ref{int:thm-ledger})
was subjected to an independent
re-derivation review: a from-scratch re-derivation of its load-bearing
implications, conducted before reading the recorded proofs and then
reconciled against them, accompanied by fresh machine checks re-implemented
from the definitions.

\paragraph{Implications.} Of the $34$ load-bearing implications of the chain,
$31$ were independently re-derived as valid. The remaining $3$ are
wording- and registration-level only: an overstated equivalence in the
density remark of the chain (the unconditional direction is one-way), an
under-registered residue inside the degree-two layer of the budget
conjecture (recorded in Section~\ref{sta:subsec-limitations}), and one
stale supporting statistic. None reverses an implication direction, and no
mathematical hole was found in any implication.

\paragraph{Machine checks, at the downgraded caliber.} Six independent
machine-check groups accompanied the review. \emph{Five of the six} closed
with zero failures (volumes: $180{,}070$ construction instances;
$550$ end-to-end pipeline instances; $876$ census graphs; $121$ types of the
placement family; $2{,}342$ buried-vertex instances; $5{,}880$ alignment
instances; plus a dedicated connected-domain cut-vertex check on $106$
instances). The sixth group closed with three of its four blocks at zero
failures and one block emitting $60$ alerts; each of the $60$ was examined
individually and all $60$ are \emph{disconnected} graphs, outside the
connected domain on which the budget conjecture is stated --- zero connected
counterexamples. The account is therefore never ``six groups with zero
failures'': the correct statement is five groups clean, and the sixth clean
up to a fully discharged false alarm caused by a missing connectivity
filter, with the connected domain covered by the dedicated check.

\paragraph{What was \emph{not} re-verified.} The same review registered its
limits, which travel with any citation of it: three combinational notes of
the chain were logic-checked but not separately machine-checked; machine
checks predating the review were not replayed; the reviewers'
derivation-before-reading process is not itself reproducible (only its
products were reconciled against the originals); the full proof
reconstruction of the all-odd branch of Theorem~\ref{int:thm-lovasz}
remains open (Section~\ref{sta:subsec-limitations}); and the domains $n\ge13$
and the general-graph domain beyond $12$ edges at $n\ge9$ were not
enumerated. (One gap registered by the review --- two of the five domains
below not freshly recomputed --- was subsequently closed by a full fresh
recomputation of all five domains, $178{,}860/178{,}860$ instances over
$16$ fields with zero differences.)

\paragraph{The five exhaustively enumerated domains.} The budget conjecture
$K$ (Definition~\ref{int:def-budget}) carries a machine account on five
exhaustively enumerated domains: the atlas of connected graphs on
$n\in\{4,6\}$ with $m\le13$ together with \texttt{geng -c 8 1:12}, and the
connected \texttt{geng} censuses \texttt{8/9/10/11/12 1:12} --- in total
$39{,}027$ graphs, $178{,}860$ even-degree point instances
(a vertex $x$ of even degree in a graph of the domain), of which $43{,}009$
have degree $d\ge4$. The boundary of the domain is part of the account: only
the sparse layers with $m\le12$ are enumerated (at $n=11$ the average degree
is $\le2.18$); the general-graph domain beyond $m\le12$ is not covered; and
odd and even $n$ are accounted separately throughout, never pooled (the
budget ceiling is $\lceil n/2\rceil$ and the analysis differs by parity).

\begin{table}[ht]
\centering\small
\caption{The five exhaustively enumerated domains of the budget conjecture.
Columns: graphs; even-degree point instances; with $d\ge4$; instances dead
at the first two insertion depths of the degree-insertion tool; instances
where the tool certifies failure; instances attaining the ceiling
$\beta=\lceil n/2\rceil$. In every domain: zero violations of
$\beta(G,x)\le\lceil n/2\rceil$, and zero mismatches of the insertion
criterion's equivalence against an independent tool-truth oracle.}
\label{com:tab-kdomains}
\begin{tabular}{lrrrrrr}
\hline
Domain & graphs & instances & $d\ge4$ & two-probe & tool-dead & tight\\
\hline
atlas $n\in\{4,6\}$, $m\le13$, $\cup$ \texttt{geng -c 8} & $2{,}933$ &
$11{,}328$ & $4{,}232$ & $306$ & $136$ & $63$\\
\texttt{geng -c 9} & $7{,}654$ & $33{,}296$ & $10{,}003$ & $227$ & $12$ &
$1{,}420$\\
\texttt{geng -c 10} & $11{,}989$ & $55{,}664$ & $13{,}471$ & $783$ & $60$ &
$72$\\
\texttt{geng -c 11} & $10{,}874$ & $52{,}126$ & $10{,}558$ & $1{,}217$ &
$0$ & $990$\\
\texttt{geng -c 12} & $5{,}577$ & $26{,}446$ & $4{,}745$ & $1{,}102$ & $0$ &
$12$\\
\hline
total & $39{,}027$ & $178{,}860$ & $43{,}009$ & $3{,}635$ & $208$ &
$2{,}557$\\
\hline
\end{tabular}
\end{table}

Two readings of Table~\ref{com:tab-kdomains} are fixed by the account.
First, \emph{tool-dead is not a counterexample}: the degree-insertion
criterion (Section~\ref{sta:subsec-limitations}) characterizes one
construction tool, and all $208$ tool-dead instances still satisfy
$\beta\le\lceil n/2\rceil$ --- indeed the zero-violation verdict of the
conjecture itself reads on all $178{,}860$ instances. Second, the witness
certificates attached to the account total $1{,}279$ ($484$ of $500$ in the
base domain plus $795$ of $800$ across the extended domains, $5$
tool-dead instances skipped by design, zero failures), each an explicit
decomposition verified three ways. The residual proof debts of this layer
are collected in Section~\ref{sta:subsec-limitations}.

\subsection{The Lean development and its trust boundaries}\label{com:subsec-lean}

The Lean development (toolchain \texttt{leanprover/lean4:4.34.0}) certifies
a precise subset of the paper's statements. We describe what is proved, how
each class of proof is trusted, and where the boundaries are; the
itemized list with file anchors is Table~\ref{sta:tab-leanitems} of
Section~\ref{sec:open}.

\paragraph{What is proved.} Two classification-line theorems are fully
formalized from first principles: for every odd $n\ge3$,
$\pn(K_{4,n})=(n+1)/2$ (with numeric corollaries $3,7,10$ at
$n=5,13,19$ following without independent gaps), and for every $k\ge1$,
$\pn(K_{2k+1})=k+1$ (via the edge-count lower bound and an explicit
Walecki-type construction; the endpoint $k=0$, where $\pn(K_1)=0$, is a
separate fact outside the theorem). Both carry the standard axiom profile
\texttt{[propext, Classical.choice, Quot.sound]} and no \texttt{sorryAx}.
On the criterion line, two declarations are formalized without
\texttt{sorryAx}: the criterion-level bound $g_{13}\le d_2+1$
(\texttt{thmA2663\_upgraded}, axioms \texttt{[propext, Quot.sound]}), where
$g_{13}$ is the criterion-layer count functional of the valuation line
(in the paper's notation, the number of ``leaks'' of $C_1$ --- points of
the shifted layer $F_{13}$ of Definition~\ref{int:def-legalval} outside
$C_1$; in the Lean sources it is defined as a count over
\texttt{List.range n} in the base module and quoted by
\texttt{ThmA2663.lean}), whose
hypotheses --- the target side has the prescribed size, and
$\mathrm{val}\le1$ --- are \emph{assumed, not proved} (it is a bound under a
hypothesis, not a feasibility-existence proof; see no-claim
\textup{(6)} of Section~\ref{int:subsec-claims}); and the case-$2$
structural discriminator (\texttt{case2\_struct\_alpha\_iff}, axioms
\texttt{[propext, Quot.sound]}), an if-and-only-if for the machine
predicate $C_1((t_1+g)\bmod N)=\texttt{true}$ under six explicit
hypotheses. Finally, the four-domain witness accounts of the termination
bound are kernel-checked at the level of the instance counts
$588/30{,}760/634{,}960/7{,}798{,}903$ and of two fully certified
five-column global ledgers:
$\mathrm{ledger}_6=(5{,}724;\,2{,}613;\,190;\,398;\,588)$
(legal pairs; witnesses; $\min\mathrm{val}=0$ bucket; $\min\mathrm{val}=1$
bucket; instances) and
$\mathrm{ledger}_8=(701{,}360;\,285{,}372;\,20{,}308;\,10{,}452;\,30{,}760)$.
The body of the termination theorem itself is proved on paper, not in Lean;
the accounts are its machine shadow.

\paragraph{The \texttt{native\_decide} boundary.} The finite certificates
(including both ledgers and the $N=10$/$N=12$ account modules) are closed by
\texttt{native\_decide}. Such theorems must be quoted as \emph{machine-checked
by kernel reduction}: they depend on the instantiation axiom of
\texttt{Lean.ofReduceBool} (the kernel-reflection axiom of the evaluation
mechanism), so they are not unproved assertions, but they do not carry the
trust level of pure term-level proofs. Three facts anchor the boundary. The
judging-layer semantics (the value functional, the board functionals, the
admission predicates) are inherited from the base module with zero
redefinition. The $N=8$ ledger is certified on two mutually pinned routes ---
whole-domain direct evaluation, and a $56$-bucket aggregation with per-bucket
certificates --- and cross-pinned against a second, independent Python
implementation, bucket by bucket. And the trust base is the honesty of the
Lean~4.34.0 kernel reflector itself.

\paragraph{The axiom audit.} The published declarations carry an axiom audit
(\texttt{\#print axioms} over the audit's probe file against the built
library). Re-running the audit against the current sources reproduces the
archived audit log entry for entry, and the six audited source files are
byte-identical to the review snapshot, so the log applies to the released
sources. The audit's complement is equally explicit: the probe covers the
declarations named in Table~\ref{sta:tab-leanitems} and the five Fan-line
declarations discussed next; anything outside the probe (the uniqueness
corollary of the Fan-type chain, the $N=10$/$N=12$ account modules, and the
auxiliary complete-graph and $K_{4,n}$ modules beyond the two named
theorems) is \emph{not} claimed fully formalized until individually probed.
A live comment-stripped scan finds exactly five \texttt{sorry} lines:
three placeholders of the Fan-type formalization chain
(\texttt{Astar.lean}, lines 2592, 2849, 2864), one \emph{refuted} research
lemma kept as a record only (\texttt{Fan.lean}, line 765 --- its statement
was disproved by a machine counterexample and is never cited; see
Remark~\ref{rem:claformal}), and one unproved research gap consumed by a
weakened variant (\texttt{Fan.lean}, line 4487). Accordingly the
Fan-type theorem chain
(Theorems~\ref{thm:claAstar} and~\ref{thm:lasep}) is cited as paper proofs
with formalization in progress, never as machine-verified
(Remark~\ref{rem:claformal}).

\paragraph{Certificate-grade example: the packing table.} The packing
residual of Section~\ref{subsubsec:criteria} carries its own Lean
certificates, which illustrate the certificate tier end to end. The main
certificate theorem checks, pointwise, all $79{,}200$ points of the table
domain ($120$ canonical orbits $\times$ $330$ demand vectors $\times$ the
two board sides): the packing predicate holds exactly when $\Sigma
d\le|S|-1$, or $\Sigma d=|S|$ with the demand outside the recorded
top-failure set. It is closed by \texttt{native\_decide} (zero
\texttt{sorryAx}, with the \texttt{ofReduceBool} disclosure attached), and
its data anchors are certified alongside it ($330$ demands; $120$ orbits;
$1{,}761$ top failures $=204+1{,}557$). Nine companion theorems over the
same table domain record: the downset property on both sides; the uniform
layer ($\Sigma d\le|S|-1$ implies feasibility on both sides); an explicit
no-threshold witness (two demands with equal $\Sigma d=4$, one feasible and
one not --- refuting any proposed threshold law of the form ``feasible
$\iff f(\Sigma d)$''); the concentration of all $1{,}761$ failures in the
top layer; the top-failure set being an upset; and the failure-orbit counts
$110$ and $2$. All are \texttt{native\_decide} certificates. Above them
sits the paper-level \emph{general} downset theorem
(\texttt{pkg\_downset}, axioms \texttt{[propext, Classical.choice,
Quot.sound]}, no \texttt{native\_decide}, no \texttt{sorryAx}), quantified
over arbitrary inputs under three slot-value hypotheses: if the slot
relation $\mathrm{slot}(d',i)+1=\mathrm{slot}(d,i)$ holds with all other
slots equal and $\mathrm{slot}(d,i)\ge1$, then feasibility of $d$ implies
feasibility of $d'$. Two boundaries travel with it: the hypotheses are
stated as slot-value relations (the subtraction form is not yet
formalized; instances are discharged through the slot relations), and the
Lean reimplementation of the packing predicate is not proven
program-level identical to the reference Python implementation --- that
identity rests on three-way machine cross-checks (the $79{,}200$-point
sweep with zero violations, bit-level reconstruction of the certificate
data, and a full Lean-side sweep).

\subsection{The key-table ledger for the covering identity}\label{com:subsec-ledger}

The machine account behind Proposition~\ref{prop:clacover} is organized as
\emph{key tables}: a key is the datum $(N,W,k,\delta_{12},\delta_{13})$ up
to rotation, and each key summarizes all instances sharing it. Across nine
domains $N\in\{6,\dots,14\}$ the key tables comprise $2{,}412{,}436$ keys
covering $98{,}602{,}396$ key-weighted admissible positions (the sum over
keys of $\binom{|W|}{d_1}$); the per-domain key counts are
$271/1{,}218/4{,}614/15{,}627/48{,}913/144{,}408/407{,}584/1{,}110{,}057/679{,}744$
for $N=6,\dots,14$. On every key: family$\setminus$prediction $=$
prediction$\setminus$family $=0$ (both inclusions of the covering identity
in key form), and the necessary-side boundary violations are $0$.

Four conventions bind every citation of these numbers; they are stated
here once and apply paper-wide.

\begin{enumerate}[label=\textup{(\arabic*)},leftmargin=2em]
\item \emph{Key-weighted versus instance-weighted.} The key tables count
each key once; instance-weighted accounts (each instance once) exist
separately for $N=8$ ($28{,}056$ instances, $221{,}480$ positions) and $N=9$
($128{,}016$ instances, $1{,}328{,}040$ positions) and are listed apart.
The two accountings are never mixed or compared.
\item \emph{Two domain conventions.} The domains $N\in\{6,\dots,13\}$ come
from exhaustive enumeration of the instance space, while $N=14$ comes from
an affine-normalized domain (orbit representatives under the normalized
action) whose equivalence to exhaustive enumeration is \emph{unproven}; the
nine-domain totals are therefore not a union under a single convention.
\item \emph{The no-$V_7$ superset reading.} All machine checks of the
covering identity are performed with the $V_7$ branch omitted
(Remark~\ref{rem:clacoverriders}), a conservative superset reading whose
zero-counterexample verdict transfers to the true face one-way. The residual
$V_7$-only mass is nonzero and precisely measured: $340$ positions at $N=8$
($1{,}266$ superset-uncovered minus $926$ true-face uncovered) and
$2{,}267{,}655$ instance-weighted at $N=14$.
\item \emph{No extrapolation.} Nothing in the ledger is extrapolated to
$N\ge15$, and the normalized domain is never read as the exhaustive domain.
\end{enumerate}

The ledger itself passed two independent controls: a re-aggregation with a
separate decoder and an independent combinatorial recomputation of the
position counts (matching the archived accounts digit for digit), and a
three-shard fresh replay with digit-identical accounting fields.

\paragraph{The independent audit.} Subsequent to, and separately from, the
verification passes above, an independent audit re-derived all twelve
supporting lemmas of the covering identity (the nine assembly lemmas and
the three micro-lemmas of Remark~\ref{rem:clacoverriders}) from the
definitional layer rather than by checking the write-up, probed each
lemma's plausible failure modes (none broke), probed each of the three
declared assembly gaps (none blocks the assembly), and re-ran the
domain-consistency check with a \emph{third} independent implementation
(exhaustive reference-form predicates, tag-level reverse direction) on
fresh key tables: four full domains ($N=6,7,8,10$; $55{,}016$ keys) plus
one $N=12$ shard ($50{,}948$ keys) --- $105{,}964$ keys in total, six
violation classes all zero, with independent aggregation matching the
headers and archived accounts digit for digit. The randomized machine
checks of the micro-lemmas were re-run in the audit pass with zero
violations ($4{,}000/1{,}109/4{,}000$ instances, after the first pass
$12{,}558/2{,}500/12{,}000$). The audit's verdict upgrades the package to
``paper candidate that has passed an independent audit'' --- which is
still not a mechanized or formalized theorem. The audit did \emph{not}
re-run, and these items are carried as archived accounts: the
single-swap falsification experiments of Section~\ref{subsubsec:criteria},
the transcription-fidelity anchor ($240$ instances, $97{,}848$ pairs), the
per-instance full-domain enumerations at $N=8,9$, and the per-instance
glue verification (the tool sources of the latter were lost with no
recoverable copy).

\appendix

\section{Computational anchors}\label{com:app-comp}

This appendix lists, for each theorem of the paper, the verification
status and the machine anchors. All computations used one pipeline
(Section~\ref{com:subsec-pipeline}); randomized steps used fixed seeds;
every number below is reproducible from the archived scripts released
with the paper. In the table, ``exact'' means computed by an exact
algorithm (dynamic program or exhaustive cover), ``witness'' means a
machine-verified decomposition, and ``census'' means all graphs in the
stated class were processed.

\begingroup\small
\begin{longtable}{p{4.6cm} p{6.2cm} p{4.2cm}}
\hline
Result & Verification & Anchors (script / certificate) \\
\hline
\endfirsthead
\hline
Result & Verification & Anchors (script / certificate) \\
\hline
\endhead
Theorem~\ref{thm:claAstar}, Theorem~\ref{thm:lasep}: witness instances &
exact $\pn\le\floor{n/2}$ on $101$ random (A1)--(A2) instances, $50$
instances with $\Delta(\EV)=5$, $15$ large witnesses; $0$ violations &
block-reload and witness-upgrade scripts\\
\hline
Theorem~\ref{thm:lasep}, Remark~\ref{rem:claCO},
Proposition~\ref{prop:claresidue}, Proposition~\ref{prop:claodddensity} &
$413$ p-anchored (A1)--(A2) instances ($13$ structured, $400$ random),
$0$ violations; exhaustive $n=5$ (all $1024$ graphs) and $n=6$ (all
$32768$ graphs) degree/triangle-layer checks behind
Proposition~\ref{prop:claodddensity} & structured-fix scripts\\
\hline
Theorem~\ref{thm:clatwohigh}, Corollary~\ref{cor:clatwohigh} (two-high
layer) & phantom-step audit $4/4$; triangle-forced exits of dirty shared
neighbours $34/34$; $34$ structured chain certificates; random two-high
census $144/144$ chains verified; exact DP cross-checks $43/43$; $0$
violations & two-high-layer scripts\\
\hline
Corollary~\ref{cor:clatwohigh}: the dirty corner & $207$ dirty-class
instances ($154$ legacy: $34$ structured $+\,120$ random; $53$ new), plus
$16$ LB1-tightness hardening runs and $60$ two-high regression chains; in
every instance a decomposition of $\tilde G$ into at most
$\floor{n/2}-1$ paths was produced (constructive refutation of the
corner); $0$ corner hits, $0$ violations & corner-narrowing scripts\\
\hline
Corollary~\ref{cor:clatwohigh}, Remark~\ref{rem:clagap}: band-free
tightness & $125$ decision-level instances of the constructed family:
all $66$ band-free instances sealed constructively (paired-decomposition
constructor output verified to partition the edge set into exactly
$\mathrm{odd}/2$ simple paths, with per-vertex endpoint assertions); all
$59$ banded instances rejected with band certificates; constructor
success $\iff$ band-free reconciled $125/125$; excess
$\mathrm{excess}\le\lfloor y/2\rfloor$ audited on these $125$ and the
$110$ earlier band-level instances of the same pipeline ($235$ in all),
zero violations & band-tightness scripts\\
\hline
Remark~\ref{rem:clagap}: the band-free statement beyond the constructed
family & $71$ decision-level instance runs: the general band-free
statement refuted by cycle-block instances certified three ways
(band-free by the criterion, untight by the decision procedure, and an
exact combinatorial witness at $\pn=\mathrm{odd}/2+1$ for the
cycle-with-one-pendant block and $\pn=\mathrm{odd}/2+2$ for the
double-triangle block); $0$ failures & relaxed-domain scripts\\
\hline
Proposition~\ref{prop:claWmbound}, Proposition~\ref{prop:claescape}:
$W(m)$ & parity certificate by hand; decompositions into exactly
$\floor{n/2}$ paths verified for $m\in\{5,7,9\}$; class memberships
asserted programmatically & witness-family scripts\\
\hline
Proposition~\ref{prop:claoscsharp}: exactness on tested $k$ & DP exact at
$k=2$; verified $(k+1)$-path witnesses for $k\in\{2,3,4,5,6,8,10\}$ &
witness-upgrade scripts\\
\hline
Theorem~\ref{thm:conk4odd} & structured exact search over $K_{4,n}$,
$n=3..31$, recovering $(n+1)/2$ independently for every $n$; closed-form
construction verified for $k=6..301$ ($n=13..603$); witness lists
machine-verified & $K_{4,n}$ scripts and witness files\\
\hline
Proposition~\ref{prop:conk4small} & witness verification (edge-exact
partition, simplicity) & small-case witness files\\
\hline
Even columns of Case~4 (Section~\ref{subsec:columns};
Theorem~\ref{thm:concapgen}, Theorem~\ref{thm:constructgen},
Theorem~\ref{thm:conringclosure}, Corollary~\ref{cor:concol6},
Observation~\ref{obs:concostwall}) & $n_2=6$: witnesses for
$K_{7,6},K_{9,6},K_{11,6},K_{13,6}$ each verified by two independent
implementations (pipeline verifier $+$ adjacency-matrix checker); ring
construction re-derived for $k=7..31$ from the base $k_0=3$, with
overlapping chains from $k_0=4,5$. $n_2=8$: bases
$K_{9,8},K_{11,8},K_{13,8}$ found at $5/7/43$ search nodes and verified
($K_{11,8},K_{13,8}$ re-derived byte-identically); closure sampled at
$60$ levels $k=9..30$; negatives on $K_{15,8},K_{17,8}$: $>14$M nodes
over all five deficit multisets, zero exhaustions. $n_2=10$: bases
$K_{11,10},K_{13,10},K_{15,10}$ found at $6/13/182$ nodes and
triple-verified; closure sampled $k=11..18$; negatives on
$K_{17,10},K_{19,10},K_{21,10}$, $\approx0.9$M nodes, zero exhaustions &
column scripts and witness files\\
\hline
Theorem~\ref{thm:conL0} (anchor) & Walecki partition checked edge by edge
for $j=2..16$; formula verified for all odd $j\le121$ & engine scripts\\
\hline
Theorem~\ref{thm:conchain} assembly & $14$ random chains
($b\in\{2..5\}$) assembled and verified; $b=4,5$ chains attain
$\sum j_i-1$ & engine scripts\\
\hline
Remark~\ref{rem:conchaincomp}(2): three $K_5$'s & the $5$-path
decomposition printed in the text is hand-checkable; contradicts the
archived exact value $6$ (re-verification requested) & this paper, project
archive\\
\hline
Conjecture~$K$ (Definition~\ref{int:def-budget}) evidence & the five
exhaustively enumerated domains of Table~\ref{com:tab-kdomains}; $0$
violations & domain scripts and certificate files\\
\hline
Proposition~\ref{prop:claincomp}: exact values & DP exact:
$\pn(K_{2,5}+2\ell)=3$, $\pn(G')=4$ (triangle plus two decorated
bridges), $n=9,11$ & witness-upgrade scripts\\
\hline
Lemma~\ref{int:lem-fan6}, both orientations (transcription and
verification note in Section~\ref{int:subsec-fanlemmas}) &
direction-corrected scan over both ordered (gains-end, budget-end)
orientations of every deleted edge: $n\le5$ complete, all $482{,}080$
ordered instances ($241{,}040$ per orientation), $0$ violations; $n=6$
complete, all $1{,}701{,}141{,}150$ ordered instances ($850{,}570{,}575$
per orientation; the fixed-orientation half exactly reproduces the
historical census) with $0$ violations & direction-scan scripts\\
\hline
\end{longtable}
\endgroup

\section{The corner analysis and the band-free tightness theorem}
\label{com:app-band}

This appendix records, with complete proofs, the theorem-level material
behind the band-free tightness statements used in
Remark~\ref{rem:clagap}, in the proof of Corollary~\ref{cor:clatwohigh},
and in open problem O2. Throughout the appendix, $G$ is a graph of the
\emph{dirty class} of Corollary~\ref{cor:clatwohigh}: $G$ is connected
and satisfies {\rm(A1)}--{\rm(A2)}, has exactly two high vertices $u$ and
$x$ with $k=\dEV(u)$ and $l=\dEV(x)$ (both odd and $\ge5$ by (A2)),
$t=t_{ux}\ge1$ common EV-neighbours $v_1,\dots,v_t$, every common
EV-neighbour of E-degree $3$, and no private EV-neighbour of E-degree $1$
on either side. We write $z_j$ for the third EV-neighbour of $v_j$, $F$
for the E-star at $u$, and $\tilde G_j=G-F-\{xv_j\}$, abbreviating
$\tilde G=\tilde G_1$. By (A1) the common EV-neighbours are pairwise
non-adjacent inside $\EV(G)$ and each $z_j$ is a \emph{far} even vertex
of $\tilde G_j$ (an even vertex outside $N_{\EV}(u)\cup N_{\EV}(x)$ that
is not a common EV-neighbour); also $\dEV(z_j)\le3$, since
$\dEV(z_j)\ge4$ would make $z_j$ a third high vertex.

\begin{definition}[corridors, terminal units, bands]\label{com:def-band}
Fix $j$ and write $\tilde G$ for $\tilde G_j$. For an even vertex $f$ of
$\tilde G$, let $Y(f)$ be its set of \emph{corridor} neighbours: the
degree-two even vertices $y$ with $N_{\tilde G}(y)=\{x,f\}$, where $N(f)$
denotes the neighbour set of $f$ in $\tilde G$. The remaining incident
edges of $f$ are its \emph{terminal units}, numbering $|N(f)|-|Y(f)|$.
Call $f$ a \emph{deadlock band} if $|Y(f)|>|N(f)|-|Y(f)|$. Let
$\mathrm{Far}$ be the set of far even vertices of $\tilde G$ and $Y$ the
total number of corridors (each corridor is a private EV-neighbour of
$x$, so $Y\le l-t$, with equality in the constructed family below).
\end{definition}

\begin{definition}[constructed family]\label{com:def-family}
$G$ belongs to the \emph{constructed family} if it lies in the dirty
class and: (i) every private vertex of E-degree $2$ (even in $G$, on
either the $u$- or the $x$-side) has its second EV-neighbour in
$\mathrm{Far}$, and no EV-edge of $\tilde G_1$ joins two far even
vertices; (ii) the odd-neighbour block of $u$ --- the substructure hanging
at $u$'s odd-degree neighbours --- is a set of $s$ pendant leaves at $u$
with $s$ odd, or a number of gadgets ($u$--$o$, $o$--$a$, $o$--$b$, with
two pendant leaves on each of $a$ and $b$), possibly mixed; and likewise
for $x$; (iii) every far vertex $f$ of $\tilde G_1$ has $\dEV(f)\le3$ and
carries at most one repair leaf, whence $\deg_{\tilde G_1}(f)\in\{2,4\}$;
and (iv) every common EV-neighbour of $G$ is adjacent to an odd positive
number of pendant leaves and to no odd-degree vertices other than these
leaves. (The band condition of Definition~\ref{com:def-band} does not
depend on $j$: no far--far EV-edge exists by (i), so deleting $xv_j$
touches no far vertex and none of its neighbours. The machine instances
of Appendix~\ref{com:app-comp} all lie in the family, with the
cross-private instances covered by the extension that follows.)
\emph{Extension.} Cross-private pairs are admitted: a $u$-private vertex
$w$ may be adjacent to an $x$-private vertex $y$, so that
$N_{\EV}(w)=\{u,y\}$ and $N_{\EV}(y)=\{x,w\}$; $y$ then acts as a
terminal corridor with $\phi(y)=w$, and $wy$ is a terminal unit of $w$.
The band-free tightness theorem below covers this extension by the
identical step (ii) of its proof; the load account and the counting
theorem are used only on the family proper, where every private even
vertex has its second EV-neighbour in $\mathrm{Far}$.
\end{definition}

\subsection{The counting gap and the counting theorem}

\begin{lemma}[counting-gap ledger]\label{com:lem-gapledger}
For any $j$, let $\mathrm{odd}:=\mathrm{odd}(\tilde G_j)$ and
$\mathrm{gap}:=\floor{n/2}-\mathrm{odd}/2$, where $n=|V(G)|$ and
$\mathrm{odd}_G$ is the number of odd-degree vertices of $G$. Then
$\mathrm{odd}=k+1+\mathrm{odd}_G$,
$n=2+t+(k-t)+(l-t)+|\mathrm{Far}|+\mathrm{odd}_G$, and
\[
\mathrm{gap}=\frac{l-t+|\mathrm{Far}|+1}{2}\ \ (n\text{ even}),
\qquad
\mathrm{gap}=\frac{l-t+|\mathrm{Far}|}{2}\ \ (n\text{ odd});
\]
in particular $\mathrm{gap}\ge1$.
\end{lemma}

\begin{proof}
The odd-degree vertices of $\tilde G_j$ are: $u$, which loses the $k$
star edges (an odd number); $x$, which loses $xv_j$; the common
EV-neighbours $v_i$ with $i\ne j$ ($t-1$ of them, each losing $uv_i$);
the private EV-neighbours of $u$ ($k-t$ of them, each losing $uw$); and
all odd-degree vertices of $G$ (untouched, since the star edges all join
even-degree vertices and $xv_j$ joins two even-degree vertices). Hence
$\mathrm{odd}=2+(t-1)+(k-t)+\mathrm{odd}_G=k+1+\mathrm{odd}_G$. The
vertex set splits into $\{u,x\}$, the common EV-neighbours, the private
EV-neighbours of $u$ and of $x$, the far even vertices (which include the
$z_i$), and the odd-degree vertices of $G$, giving the displayed count.
For odd $n$: the remaining terms $2$, $k+l$ and $\mathrm{odd}_G$ are even,
so $n$ is odd exactly when $t+|\mathrm{Far}|$ is odd, which makes
$l-t+|\mathrm{Far}|$ even and the second quotient an integer. Finally
$l\ge t$ and $|\mathrm{Far}|\ge1$ (as $z_j\in\mathrm{Far}$): for even $n$
the numerator is $\ge2$; for odd $n$ it is a positive even number, hence
$\ge2$. So $\mathrm{gap}\ge1$.
\end{proof}

\begin{lemma}[load account]\label{com:lem-load}
In the family proper (the constructed family without cross-private
pairs), $\sum_{f\in\mathrm{Far}}\dEV(f)=k+l-t$ (E-degrees read in $G$).
Consequently $\dEV(f)\le3$ for every $f\in\mathrm{Far}$, and
$|\mathrm{Far}|\ge2$.
\end{lemma}

\begin{proof}
Each common EV-neighbour $v_i$ has EV-neighbours $u$, $x$ and
$z_i\in\mathrm{Far}$, contributing $t$ EV-edges with one end in
$\mathrm{Far}$; each private even vertex (of E-degree $2$ in $G$, on
either side) has its second EV-neighbour in $\mathrm{Far}$ by
Definition~\ref{com:def-family}(i), contributing the $k-t$ edges at the
private EV-neighbours of $u$ and the $l-t$ edges at the corridors of $x$.
By (i) there are no EV-edges between far vertices, and odd-degree
vertices carry no EV-edges; so every EV-edge counted has exactly one end
in $\mathrm{Far}$, and summing the E-degrees over $\mathrm{Far}$ gives
$t+(k-t)+(l-t)=k+l-t$. If some $f$ had $\dEV(f)\ge4$, then $f$ would be a
third high vertex, impossible; so all E-degrees in $\mathrm{Far}$ are
$\le3$. If $|\mathrm{Far}|=1$, then $k+l-t\le3$, so
$t\ge k+l-3\ge5+5-3=7>\min(k,l)\ge t$, a contradiction; hence
$|\mathrm{Far}|\ge2$.
\end{proof}

\begin{proposition}[the counting theorem: small gap forces band-free]
\label{com:prop-gapnoband}
In the family proper (the constructed family without cross-private
pairs), if $\mathrm{gap}\le2$, then $\tilde G_j$ has no deadlock band,
for every $j$.
\end{proposition}

\begin{proof}
By Definition~\ref{com:def-family} the band condition is
$j$-independent, so we argue in $\tilde G_1$. Suppose $f$ is a band. Then
at least two corridors hang at $f$, and every corridor is a private
EV-neighbour of $x$; write $y:=l-t\ge Y\ge2$.
If $n$ is even, Lemma~\ref{com:lem-gapledger} gives
$y+|\mathrm{Far}|\le 2\,\mathrm{gap}-1\le3$, while $y\ge2$ and
$|\mathrm{Far}|\ge2$ (Lemma~\ref{com:lem-load}) give
$y+|\mathrm{Far}|\ge4$ --- a contradiction. If $n$ is odd,
$\mathrm{gap}\le2$ gives $y+|\mathrm{Far}|\le4$, so
$y\le4-|\mathrm{Far}|\le2$; equality throughout forces $|\mathrm{Far}|=2$
and $y=2$. The load account gives $k+l-t=k+y=k+2\le3|\mathrm{Far}|=6$, so
$k\le4$, contradicting $k$ odd and $\ge5$. So $y=2$ is impossible, and
the bound $y\le2$ then leaves $y\le1$, i.e.\ $l-t\le1$: no band exists (a
band needs two corridors at one far vertex).
\end{proof}

\subsection{The pairing framework and the band-free tightness theorem}

\begin{proposition}[pairing schemes produce trail decompositions]
\label{com:prop-pairframe}
Let $H$ be a finite simple graph with a local structure on the edges at
each vertex: if $\deg_H(v)$ is even, its incident edges are partitioned
into pairs; if $\deg_H(v)$ is odd, exactly one incident edge is marked
\emph{free} and the rest are paired. Then $E(H)$ decomposes into maximal
alternating trails --- starting from a free mark, at each internal visit
leaving through the mate of the arriving edge --- and every edge lies on
exactly one trail. Each trail either has two free ends or is closed; the
open trails number $\mathrm{odd}(H)/2$. If moreover no trail repeats a
vertex, then $H$ has a path decomposition $D$ with
$|D|=\mathrm{odd}(H)/2$, every odd-degree vertex an endpoint of exactly
one path of $D$ and no even-degree vertex an endpoint; in particular
$\pn(H)=\mathrm{odd}(H)/2$.
\end{proposition}

\begin{proof}
At each vertex the successor of an arriving edge is uniquely determined
by the pairing, so the trails are edge-disjoint and exhaust $E(H)$. A
free mark can occur only at a trail end (a free edge carries at most one
free mark at each of its ends); an odd vertex has exactly one free mark
and is therefore the end of exactly one trail; an even vertex has none
and is internal wherever it occurs. The number of open trails is the
number of free marks divided by two, i.e.\ $\mathrm{odd}(H)/2$. If all
trails are paths, they form a decomposition $D$ with $|D|=\mathrm{odd}/2$
and endpoint set exactly the odd vertices; the odd-degree lower bound of
Lemma~\ref{int:lem-LB1}(i) gives $\pn(H)\ge\mathrm{odd}/2$, so the value
is exact.
\end{proof}

\begin{lemma}[odd-neighbour blocks close by themselves]\label{com:lem-blocks}
In the constructed family, the odd-neighbour block of $c\in\{u,x\}$ can
be structured as Proposition~\ref{com:prop-pairframe} requires, with all
resulting trails staying inside the block, as follows. A leaf block of
$s$ leaves ($s$ odd): mark $c\ell_1$ free at $c$ (the single-edge path
$\ell_1$--$c$) and pair $(c\ell_{2i-1},c\ell_{2i})$ for $2\le 2i\le s$,
giving $(s+1)/2$ paths for the $s+1$ odd vertices involved. A gadget:
pair $c$--$o$ free at $c$ (path $c$--$o$), $(o$--$a,o$--$b)$ at $o$,
$(a\ell_1,a o)$ at $a$ with $a\ell_2$ free, and $(b o,b\ell_3)$ at $b$
with $b\ell_4$ free; the resulting paths $c$--$o$,
$\ell_1$--$a$--$o$--$b$--$\ell_3$, $\ell_2$--$a$, $\ell_4$--$b$ have
endpoint multiset covering $\{c,o,a,b,\ell_1,\dots,\ell_4\}$ exactly once
--- $8$ odd vertices, $4$ paths. A mixed block (gadgets and leaves at the
same $c$): the block edges at $c$ number odd in total, so exactly one is
free (chosen canonically, say the first), the rest pair arbitrarily among
leaf and gadget edges, and a mixed pair $(c\ell,c o)$ forms the legal
path $\ell$--$c$--$o$. All trails end at free marks inside the block, so
the block never interacts with the rest of the scheme.
\end{lemma}

\begin{proof}
This is a verification of the pairing tables just displayed: within a
gadget the pairs at $o,a,b$ are fixed as shown; leaves pair among the
edges at $c$; and in a mixed block the total number of block edges at $c$
is $s+g$ ($s$ leaves, $g$ gadgets, each gadget contributing a single
attachment edge $co$), which is odd because $c$ has even degree in $G$
and an odd number of EV-neighbours, so the free/paired structure exists;
the endpoint count of each displayed path list matches the odd vertices
it involves, each exactly once. No trail leaves the block, since free
marks sit at block vertices only.
\end{proof}

\begin{lemma}[forced fusions]\label{com:lem-fuse}
In any decomposition of $\tilde G$ as in the last sentence of
Proposition~\ref{com:prop-pairframe} (a \emph{paired decomposition}), the
edges at every even vertex are traversed in internal pairs. Hence a
corridor $y$ (degree $2$) forces its two edges $xy$ and $y\phi(y)$ into
one path; and the unit at $v_1$ is forced: $v_1z_1$ pairs with
$v_1\ell_1$, where $\ell_1$ is the leaf at $v_1$ (any further leaves at
$v_1$ pair among themselves, their number being even by
Definition~\ref{com:def-family}(iv)).
\end{lemma}

\begin{proof}
At an even vertex $v$ of a paired decomposition, $D(v)=0$, so no path has
$v$ as an endpoint; each visit is internal and uses exactly two edges, so
the incident edges split into internal pairs by path. A degree-two vertex
therefore has its whole edge set as a single pair on a single path.
\end{proof}

\begin{proposition}[injection at far vertices]\label{com:prop-inject}
Fix a far vertex $f$ of $\tilde G$ with corridor set $Y(f)$, $c(f):=|Y(f)|$,
and terminal units $T(f)$, $|T(f)|=\deg(f)-c(f)$. The edges at $f$ can be
paired so that every corridor edge $fy$ is paired with a \emph{distinct}
terminal unit if and only if $f$ is not a band, i.e.\ $c(f)\le|T(f)|$ (the
pairing at $f$ is otherwise unconstrained, so this is a matching condition
in a complete bipartite graph between $Y(f)$ and $T(f)$). In the
constructed family $c(f)\le2$.
\end{proposition}

\begin{proof}
Only the bound remains: by Definition~\ref{com:def-family}(iii)
$\dEV(f)\le3$ and $f$ carries at most one repair leaf, so
$\deg(f)\in\{2,4\}$, and $c(f)\le\deg(f)/2\le2$.
\end{proof}

\begin{lemma}[sorted offset matching at $x$]\label{com:lem-offset}
In the constructed family, $x$ is odd in $\tilde G$, with exactly one
free edge; its non-block edges are the $m$ corridor edges (each labelled
by the far vertex it hangs at, with multiplicity $c(f)\le2$ per label; a
cross-private corridor counts as a label of its own, since no two of them
hang at the same $w$) and the $s=t-1$ bridge edges $xv_i$ ($i\ge2$), with
$m+s=l-1$ even. There is a pairing of these edges in which no two paired
corridor edges carry the same label:
if $m$ is even, sort the corridors $p_1,\dots,p_m$ by (label, index) and
pair $(p_i,p_{i+m/2})$ for $i=1,\dots,m/2$, pairing the bridges among
themselves --- if $p_i,p_{i+m/2}$ shared the label $f$, the label $f$
would occupy at least $m/2+1$ of the sorted positions, so
$c(f)\ge m/2+1\ge3$ for $m\ge4$, contradicting
Proposition~\ref{com:prop-inject};
if $m=2$ with equal labels, pair each of the two corridors with a bridge
and the remaining bridges among themselves ($s-2$ is even since
$s\equiv m\pmod2$); the failing configuration $m=2$, equal labels, $s=0$
is unreachable, since $s=0$ forces $t=1$ and $m=l-1$, i.e.\ $l=3$,
contradicting $l\ge5$;
if $m$ is odd (possible only when $t$ is even), extract one corridor of a
label of maximum multiplicity together with one bridge into a mixed pair,
and apply the even case to the rest.
\end{lemma}

\begin{proof}
The counting arguments displayed are the whole proof: for the even case,
if both $p_i$ and $p_{i+m/2}$ carried label $f$, then all positions from
$i$ to $i+m/2$ would carry label $f$ (the sort is by label first), giving
$m/2+1$ corridors at $f$. For $m=2$ with equal labels and $s\ge2$, the
corridor--bridge pairs leave $s-2\ge0$ bridges, even in number. For odd
$m$, extracting one corridor of a maximum-multiplicity label and one
bridge leaves an even number of each, and the extracted pair mixes a
corridor with a bridge, so the no-equal-label condition is preserved.
\end{proof}

\begin{theorem}[band-free tightness in the constructed family]
\label{com:thm-bandfree}
Let $G$ belong to the constructed family and fix $j$. If $\tilde G_j$ has
no deadlock band, then $\tilde G_j$ admits a paired decomposition;
consequently $\pn(\tilde G_j)=\mathrm{odd}(\tilde G_j)/2$.
\end{theorem}

\begin{proof}
We define the local structure of Proposition~\ref{com:prop-pairframe} on
$\tilde G_j$.
(i) The odd-neighbour blocks of $u$ and of $x$ are structured as in
Lemma~\ref{com:lem-blocks}, $x$'s free edge being the canonical first
edge of its block.
(ii) The degree-two even vertices are fused as in Lemma~\ref{com:lem-fuse}:
each corridor $y$ pairs $xy$ with $y\phi(y)$, where
$\phi(y)\in\mathrm{Far}$ is the far vertex at its other end (for the
cross-private \emph{extension} of Definition~\ref{com:def-family} the
same fusion applies with $\phi(y)=w$ a $u$-private vertex, $wy$ being a
terminal unit of $w$, so the extension is covered by the identical
argument); the unit at $v_j$ pairs $v_jz_j$ with $v_j\ell_j$. A shared
vertex $v_i$, $i\ne j$, takes the pair $(v_ix,v_i\ell_i)$, marks $v_iz_i$
free, and self-pairs its remaining leaf edges.
(iii) At each far vertex $f$, Proposition~\ref{com:prop-inject} ---
applicable because $\tilde G_j$ is band-free --- pairs every corridor
edge $fy$ with a distinct terminal unit; the remaining terminal units
number $\deg(f)-2c(f)$, which is even, and pair among themselves.
(iv) At $x$, the free edge is the block edge of (i), and the remaining
edges pair according to Lemma~\ref{com:lem-offset}.
(v) Every degree-one vertex carries the free mark at its leaf end. The
free-mark audit is immediate: each odd vertex of $\tilde G_j$ carries
exactly one free mark and each even vertex none, by construction.

\emph{No trail closes.} Consider any trail of the resulting scheme. At a
far vertex a corridor end is always paired with a terminal unit, so a
trail contains at most two corridors and, having left a far vertex, never
returns to it; at $x$ each trail passes exactly once (a single pair); the
fusions at $v_j$ and the bridges at the $v_i$ are one-way terminations.
Hence every trail has one of the forms
\[
\text{terminal}--f--y--x--y'--f'--\text{terminal}',\qquad
\text{terminal}--f--y--x--v_i--\ell_i,
\]
\[
\ell_i--v_i--x--v_{i'}--\ell_{i'},\qquad
\text{terminal}--f--\text{terminal}',
\]
where in the first form $y$ and $y'$ are corridors of \emph{distinct} far
vertices $f\ne f'$, by Lemma~\ref{com:lem-offset}; a ``terminal'' means a
terminal unit followed to its free end --- the whole segment from $f$
along a repair-leaf edge, along a hook edge to a $u$-private vertex $w$,
or along the $z$-role edge $v_iz_i$ read from $z_i$ to its free mark at
$v_i$. When the corridor $y'$ of the first form is cross-private, its
right-hand segment $y'$--$f'$--terminal$'$ is replaced by $y'$--$w$: the
trail ends at the free mark of $w$, an odd vertex of $\tilde G$ (so no
closed trail arises). In each form the vertices are pairwise distinct, so
every trail is a simple path. Proposition~\ref{com:prop-pairframe} yields
a paired decomposition, and the value of $\pn$ follows.
\end{proof}

\begin{lemma}[a band blocks paired decompositions]\label{com:lem-bandblocks}
If $\tilde G_j$ has a deadlock band $f$, then $\tilde G_j$ admits no
paired decomposition; consequently
$\pn(\tilde G_j)\ge\mathrm{odd}(\tilde G_j)/2+1$.
\end{lemma}

\begin{proof}
Two corridors $y_1\ne y_2$ of the same band $f$ cannot lie on one path of
any decomposition: if a path contained both $fy_1$ and $fy_2$, then
$y_1$, being of degree two and internal (even vertices are never
endpoints in a paired decomposition), would force $xy_1$ onto the same
path, likewise $xy_2$, so the path would contain all four edges of the
cycle $x$--$y_1$--$f$--$y_2$--$x$ --- impossible for a simple path. Now in
a paired decomposition the $\deg(f)$ edges at $f$ split into
$\deg(f)/2$ internal pairs. The edge $fy$ of a corridor $y\in Y(f)$
cannot pair with another corridor edge $fy'$ (previous sentence), so it
pairs with an edge at a terminal unit; the resulting assignment
$y\mapsto$ its paired edge is injective, whence
$c(f)\le|N(f)|-c(f)$, contradicting the band condition. So no paired
decomposition exists, i.e.\ no decomposition with $\mathrm{odd}/2$ paths
and endpoint set the odd vertices; with fewer than $\mathrm{odd}/2$ paths
excluded by Lemma~\ref{int:lem-LB1}(i), this gives
$\pn\ge\mathrm{odd}/2+1$.
\end{proof}

\begin{proposition}[one band costs exactly one path]\label{com:prop-patch}
Let $G$ belong to the constructed family and suppose $\tilde G_j$ carries
exactly one deadlock band $f$. Then
$\pn(\tilde G_j)=\mathrm{odd}(\tilde G_j)/2+1$.
\end{proposition}

\begin{proof}
For a single band the pair $(c(f),|T(f)|)$ is $(2,0)$ or $(3,1)$: the
band condition $c(f)>|T(f)|$ with $c(f)+|T(f)|=\deg(f)\in\{2,4\}$ and
$c(f)\ge2$ leaves exactly these two cases, and in both the surplus of
corridor ends over terminal units has size $|S(f)|=2$; write
$S(f)=\{a,a'\}$. Modify the scheme of Theorem~\ref{com:thm-bandfree} in
three places only:
(1) at $f$, self-pair the surplus edges $(fa,fa')$ instead of injecting
them;
(2) at $x$, spend the free edge on $xa$ rather than on a block edge,
marking $xa$ free at both ends (the single-edge path $x$--$a$);
(3) the remaining injections at $f$ continue as before (in the case
$(3,1)$ the one remaining corridor pairs with the one terminal unit);
every other edge follows Theorem~\ref{com:thm-bandfree}, whose offset
matching at $x$ keeps distinct labels apart automatically, including among
the remaining corridors of the band.
The resulting trails are $a$--$f$--$a'$--$x$--\dots--terminal (with $a'$
fused, $x$ passed once, $f$ passed once) and the single edge $x$--$a$;
all are simple, and the $4$-cycle $x$--$a$--$f$--$a'$--$x$ is cut at both
$x$ and $a$ by the free marks of $xa$, so it does not arise as a trail.
The endpoint count: $a$ gains two ends (one on the $x$ side, one on the
$f$ side) and every other vertex is unchanged, so the number of path-ends
is $\mathrm{odd}+2$ and the number of paths is $\mathrm{odd}/2+1$. The
reverse inequality is Lemma~\ref{com:lem-bandblocks}.
\end{proof}

\begin{proposition}[several bands]\label{com:prop-excb}
If $\tilde G_j$ carries $B$ deadlock bands, then
$\pn(\tilde G_j)\le\mathrm{odd}(\tilde G_j)/2+2B-1$; moreover
$B\le\lfloor Y/2\rfloor$, where $Y$ is the total number of corridors.
\end{proposition}

\begin{proof}
Start from the scheme of Theorem~\ref{com:thm-bandfree} and repair the
bands one at a time. The first band is repaired as in
Proposition~\ref{com:prop-patch} (cost $+1$, spending $x$'s free edge and
one ordinary matching partner). Each further band is repaired by the
double-endpoint modification: self-pair the band's surplus pair at its
far vertex $f$ and self-pair two $x$-edges at $x$, freeing four ends; the
two resulting trails $a$--$x$--$a'$ and $a$--$f$--$a'$ add two ends, that
is, one path, per band --- an exchange costing two endpoint pairs for one
path, relying on no matching partner. Hence
$p\le\mathrm{odd}/2+1+2(B-1)=\mathrm{odd}/2+2B-1$. Finally, distinct
bands have disjoint corridor sets (each corridor hangs at exactly one far
vertex) and each band has $c\ge2$, so
$Y=\sum_{f\text{ band}}c(f)\ge2B$.
\end{proof}

\begin{corollary}[bands are exactly the obstructions]\label{com:cor-bandiff}
In the constructed family: $\tilde G_j$ has a deadlock band if and only
if $\pn(\tilde G_j)>\mathrm{odd}(\tilde G_j)/2$. In particular, on the
family proper (the constructed family without cross-private pairs), if
the counting gap is at most $2$, then
$\pn(\tilde G_j)=\mathrm{odd}(\tilde G_j)/2=\floor{n/2}-\mathrm{gap}$ for
every $j$; the cross-private instances at gap at most $2$ are covered at
probe level by the machine certificates of Appendix~\ref{com:app-comp}.
\end{corollary}

\begin{proof}
One direction is Lemma~\ref{com:lem-bandblocks}, the other
Theorem~\ref{com:thm-bandfree}. If $\mathrm{gap}\le2$,
Proposition~\ref{com:prop-gapnoband} gives band-free for every $j$ on the
family proper, and Theorem~\ref{com:thm-bandfree} with
Lemma~\ref{com:lem-gapledger} gives
$\pn=\mathrm{odd}/2=\floor{n/2}-\mathrm{gap}$.
\end{proof}

\subsection{The five escape routes at a common EV-neighbour}

\begin{theorem}[escape routes]\label{com:thm-escape}
Let $G$ be a dirty-class graph, fix a common EV-neighbour $v_1$ with
third EV-neighbour $z_1$, write $B=\floor{n/2}$ and $\tilde G=\tilde G_1$,
and let $Y$ be the set of private EV-neighbours of $x$. If any of the
following holds, then $\pn(G)\le B$:
\begin{enumerate}[label=\rm(\alph*),leftmargin=2em]
\item $\pn(\tilde G)\le B-1$;
\item $\tilde G$ has a decomposition $\tilde D$ with $|\tilde D|\le B$ and
$\tilde D(x)\ge3$;
\item $\tilde G$ has a decomposition $\tilde D$ with $|\tilde D|\le B$,
$\tilde D(x)=1$ and $\tilde D(z_1)\ge1$;
\item $\tilde G$ has a decomposition $\tilde D$ with $|\tilde D|\le B$,
$\tilde D(v_1)\ge2$ and $\tilde D(v_1)>\#\{y\in Y:\tilde D(y)=0\}$;
\item $\tilde G$ has at least two components of odd order.
\end{enumerate}
\end{theorem}

\begin{proof}
The reinsertion skeleton is that of the proof of
Theorem~\ref{thm:clatwohigh}. (s1) Reinsert $xv_1$ towards $v_1$ by
Lemma~\ref{int:lem-fan6}; the condition to check is
$\tilde D(x)>\#\{w\in N_{\tilde G}(v_1):\tilde D(w)=0\}$. Among the
neighbours of $v_1$ in $\tilde G$, the only even-degree vertex is $z_1$:
$x$ is odd in $\tilde G$ (it lost $xv_1$), and every other neighbour of
$v_1$ is an odd-degree vertex of $G$ (the edge $uv_1$ is deleted, and
$v_1$ has E-degree $3$ with EV-neighbours $u,x,z_1$). So the passing
count at $v_1$ is at most $1$. (batch) Lemma~\ref{int:lem-fan7} then
absorbs the whole star $F$ towards $u$, returning $A\subseteq F$ with
$|A|\ge\ceil{k/2}$; its hypothesis holds read over $G$ because $v_1$ is
now active and every other star vertex is odd in $\tilde G$, hence has
$\tilde D\ge1$ by Lemma~\ref{int:lem-LB1}(ii). (L) Each edge of
$F\setminus A$, $|F\setminus A|=\floor{k/2}$, is reinserted towards its
star vertex as in the proof of Theorem~\ref{thm:claAstar}: dead
candidates are confined to $N_{\EV}(v)\setminus\{u\}$, of size at most
$\dEV(v)-1\le2$ by (A2), against a budget of at least $3$ at the last
step.

(b) $\tilde D(x)\ge3>1\ge$ passing$(v_1)$, so (s1) applies and activates
$v_1$; (batch) and (L) follow.

(a) Let $D^*$ be optimal with $|D^*|\le B-1$. If $D^*(x)\ge3$, route (b)
applies. Otherwise $D^*(x)=1$ (it is odd by Lemma~\ref{int:lem-LB1}(iii),
$x$ being odd in $\tilde G$). Now
$\deg_{\tilde G}(x)=\deg_G(x)-1\ge l+1-1\ge5$ (as $x$ is high,
$\deg_G(x)\ge l+1$), so $x$ is internal to some path of $D^*$; cutting
that path at $x$ raises $|D^*|$ by one --- still at most $B$ --- and
raises the endpoint count at $x$ to $3$: route (b).

(c) $\tilde D(z_1)\ge1$ makes the passing count at $v_1$ equal to $0$, so
(s1) applies unconditionally; (batch) and (L) as in (b).

(d) Reverse the order of (s1): reinsert $xv_1$ \emph{towards $x$} first.
The Lemma~\ref{int:lem-fan6} condition is now applied at $v_1$ as the
donor: $\tilde D(v_1)>\#\{w\in N_{\tilde G}(x):\tilde D(w)=0\}$, and the
dead candidates at $x$ are precisely the dead private EV-neighbours of
$x$ (the other neighbours of $x$, namely $v_2,\dots,v_t$ and the
odd-degree neighbours, all have $\tilde D\ge1$); hypothesis (d) is
exactly this inequality. Afterwards $v_1$ is still active
($\tilde D(v_1)-1\ge1$), so Lemma~\ref{int:lem-fan7} absorbs the
\emph{full} star $F$, including $uv_1$, towards $u$; and the (L) step
runs as above, $u$'s budget being untouched by the first step. The
endpoint count at $x$ ends at $\tilde D(x)+1$, which is even, consistent
with $x$ being even in $G$; all $k+1$ deleted edges are restored and
$|D|=|\tilde D|\le B$.

(e) Every component of $\tilde G$ lies in the class of
Theorem~\ref{int:thm-fan}: it is triangle-free (an induced subgraph of
$\EV(G)$) and every E-degree is at most $3$ (the private EV-neighbours of
$x$ by (A2) at $x$, the far vertices because $u$ and $x$ are the only
high vertices, $v_1$ because $\dEV(v_1)=3$). Component-wise
Theorem~\ref{int:thm-fan} with Lemma~\ref{int:lem-sum} gives
$\pn(\tilde G)\le\sum_i\floor{n_i/2}=(n-c_{\mathrm{odd}})/2$, where
$c_{\mathrm{odd}}\ge2$ is the number of odd-order components (it is even
when $n$ is even, and odd and $\ge3$ when $n$ is odd, giving
$(n-3)/2=(n-1)/2-1=B-1$). In both cases $\pn(\tilde G)\le B-1$: route
(a).
\end{proof}

\begin{corollary}[the corner, formally]\label{com:cor-corner}
Let $G$ be a dirty-class graph with $\pn(G)>\floor{n/2}$. Then for
\emph{every} common EV-neighbour $v_j$ (with third EV-neighbour $z_j$ and
$\tilde G_j=G-F-\{xv_j\}$), all of the following hold:
$\pn(\tilde G_j)=\floor{n/2}$; $\beta(\tilde G_j,x)=\floor{n/2}+1$; $z_j$
is buried (endpoint count $0$) in every optimal decomposition of
$\tilde G_j$; in every optimal decomposition $\tilde D$, either
$\tilde D(v_j)=0$, or $\tilde D(v_j)\ge2$ and at least $\tilde D(v_j)$
private EV-neighbours of $x$ have $\tilde D=0$; and $\tilde G_j$ has at
most one component of odd order.
\end{corollary}

\begin{proof}
This is the contrapositive of Theorem~\ref{com:thm-escape}, route by
route, for each $j$. First $\pn(\tilde G_j)\le\floor{n/2}$ always holds:
every E-degree in $\tilde G_j$ is at most $3$ (as in the proof of route
(e)), so component-wise Theorem~\ref{int:thm-fan} with
Lemma~\ref{int:lem-sum} applies. Route (a) failing then gives
$\pn(\tilde G_j)=\floor{n/2}$. Route (b) failing, with endpoint counts at
$x$ odd (Lemma~\ref{int:lem-LB1}(iii), $x$ being odd in $\tilde G_j$),
gives $\tilde D(x)=1$ in every optimal decomposition; cutting a path of
an optimal decomposition at the internal vertex $x$ (legal since
$\deg_{\tilde G_j}(x)=\deg_G(x)-1\ge5$) yields $\floor{n/2}+1$ paths with
endpoint count $3$ at $x$, so $\beta(\tilde G_j,x)=\floor{n/2}+1$. Route
(c) failing gives $\tilde D(z_j)=0$ in every optimal decomposition.
Route (d) failing gives the fourth condition; when $l-t\le1$ it
degenerates to $\tilde D(v_j)=0$, since at most one private EV-neighbour
of $x$ can be dead while $\tilde D(v_j)\ge2$ would need at least two.
Route (e) failing gives at most one odd-order component.

(A note on conventions: the $\beta$ here is the $D(x)\ge2$ variant of
Definition~\ref{int:def-budget}. The $D(x)\ge3$ variant takes the same
value $\floor{n/2}+1$: its lower bound comes from route (b) failing
exactly as above, and its upper bound from the same cutting construction,
which raises the endpoint count at $x$ to $3$.)
\end{proof}

% end sec-computation.tex
